% =====================================================================
%  Vol I Climax Theorem: the two equations at the Part II -> Part III hinge
%  Inscription target: chapters/theory/ or chapters/connections/
%  Produced by Wave 17 opus_10 attack-heal cycle, 2026-04-24.
%
%  The two equations
%   (CL-1)   d_bar  =  KZ^*(nabla_Arnold)
%   (CL-2)   K(A)   =  - c_ghost(BRST(A))
%  together collapse Theorems A, B, C, D, H into corollaries.
%  Attack-heal cycles exhibit (a) the derivation of each theorem,
%  (b) the Trinity identification K_E = K_c = K_g, (c) concrete
%  verification on Heisenberg, V_k(sl_2), Virasoro, beta-gamma.
%
%  Dependencies assumed in ambient file:
%      \providecommand{\fg}{\mathfrak{g}}
%      \providecommand{\ChirAlg}{\operatorname{ChirAlg}}
%      \DeclareMathOperator{\Conf}{Conf}
%      \providecommand{\Ran}{\operatorname{Ran}}
%  Theorems, lemmas, claim-status tags standard (see main.tex preamble).
% =====================================================================

\chapter[The two equations]{The two equations: \(d_{\mathrm{bar}} = \mathrm{KZ}^{*}\,\nabla_{\mathrm{Arnold}}\) and \(K(A) = -c_{\mathrm{ghost}}\,\mathrm{BRST}(A)\)}
\label{ch:climax-two-equations}

\providecommand{\KZ}{\mathrm{KZ}}
\providecommand{\Arn}{\nabla_{\mathrm{Arnold}}}
\providecommand{\dbar}{d_{\mathrm{bar}}}
\providecommand{\cghost}{c_{\mathrm{ghost}}}
\providecommand{\BRST}{\mathrm{BRST}}
\providecommand{\Bord}{B^{\mathrm{ord}}}
\providecommand{\ChirAlgInf}{\ChirAlg^{E_{\infty}}}
\providecommand{\HH}{\mathrm{HH}}
\providecommand{\Vir}{\mathrm{Vir}}
\providecommand{\Heis}{\mathcal{H}}
\providecommand{\DS}{\mathrm{DS}}
\providecommand{\KSDual}{\mathrm{KSDual}}
\providecommand{\Deldual}{\mathbb{D}}
\providecommand{\Hom}{\mathrm{Hom}}

The chiral bar complex carries a differential and the chiral algebra
carries a characteristic. Neither arrives from outside: Arnold's 1969
three-term identity on configurations of points on a smooth curve
supplies the differential, and the Friedan--Martinec--Shenker 1986
ghost-charge formula supplies the characteristic. Everything else in
Volume I --- the bar--cobar adjunction, the Koszul inversion, the
complementarity ceiling, the genus-universality, and the Hochschild
concentration --- follows from these two inputs read through one
functor. The chapter inscribes the unification.

\section{Setup}
\label{sec:climax-setup}

Fix a smooth projective curve \(X\) over \(\C\). Let \(\ChirAlgInf\)
denote the category of \(E_{\infty}\)-chiral algebras on \(X\) in the
sense of Beilinson--Drinfeld: quasi-coherent sheaves of augmented
unital factorisation \(D\)-modules with chiral multiplication
\(
  \mu \colon j_{!}(A \boxtimes A)|_{U} \to \Delta_{!}\,A
\)
on the open complement \(U\) of the diagonal, satisfying associativity,
commutativity, and the unit axiom. For \(A \in \ChirAlgInf\) augmented
by \(\epsilon\colon A \to \mathbb{1}\), the reduced part is
\(\bar A = \ker \epsilon\); the ordered chiral bar complex is
\[
  \Bord(A) \;=\; T^{c}(s^{-1} \bar A) \;=\;
  \bigoplus_{n \geq 0} (s^{-1} \bar A)^{\otimes n},
\]
realised geometrically as the factorisation module on
\(\Conf_{n}^{\mathrm{ord}}(X)\) produced by the Fulton--MacPherson
residue construction of Part~I. The differential \(\dbar\) is the sum
of contraction maps along codimension-one collision strata.

On the configuration side, write \(\Conf_{n}^{\mathrm{ord}}(X)\) for
the ordered configuration space and
\(\eta_{ij} = d\log(z_{i} - z_{j})\) for the canonical logarithmic
one-forms on the open part
\(\Conf_{n}^{\mathrm{ord}}(X) \cap \mathbb{A}^{n}\). Let
\(\mathfrak{t}_{n}\) be the infinitesimal braid Lie algebra on
generators \(t_{ij} = t_{ji}\), \(1 \leq i \neq j \leq n\), modulo
\(
  [t_{ij}, t_{kl}] = 0
\) when \(\{i,j\} \cap \{k,l\} = \emptyset\),
\(
  [t_{ij}, t_{ik} + t_{jk}] = 0
\) for distinct \(i,j,k\). Arnold's universal KZ connection is
\[
  \Arn \;=\; d \;-\; \sum_{i<j} t_{ij} \cdot \eta_{ij}
  \quad\text{on } \Conf_{n}^{\mathrm{ord}}(X) \otimes U(\mathfrak{t}_{n}),
\]
with flatness \(\Arn^{2} = 0\) equivalent to the three-term identity
\(\eta_{ij}\wedge\eta_{jk} + \eta_{jk}\wedge\eta_{ki} + \eta_{ki}\wedge\eta_{ij} = 0\)
plus the infinitesimal braid relations.

\begin{definition}[\(bc\)-ghost central charge]
\label{def:bc-charge-climax}
For \(\lambda \in \tfrac12\Z_{>0}\), the spin-\(\lambda\) fermionic
\(bc(\lambda)\) system \((b(z),c(z))\) with
\(b(z)c(w) \sim (z-w)^{-1}\) carries Virasoro central charge
\(c_{bc(\lambda)} = -2(6\lambda^{2} - 6\lambda + 1)\). The bosonic
\(\beta\gamma(\lambda)\) system is obtained by reversing the fermionic
parity, with \(c_{\beta\gamma(\lambda)} = +2(6\lambda^{2}-6\lambda+1)\).
\end{definition}

\begin{definition}[Quasi-free BRST resolution]
\label{def:brst-resolution-climax}
A \emph{quasi-free BRST resolution} of \(A \in \ChirAlgInf\) is a
quasi-isomorphism of chiral cdgas
\(
  (C_{\bullet}, Q) \xrightarrow{\sim} A
\)
with \(C_{\bullet}\) a free graded-commutative chiral algebra on
generators \((\lambda_{\alpha}, \varepsilon_{\alpha})\) of conformal
weights \(\lambda_{\alpha}\) and \(\Z/2\)-grades
\(\varepsilon_{\alpha} \in \{0,1\}\), and \(Q\) an odd-degree derivation
of square zero. The total ghost central charge is
\(
  \cghost(C_{\bullet})
    = \sum_{\alpha} (-1)^{\varepsilon_{\alpha}} \cdot 2(6\lambda_{\alpha}^{2} - 6\lambda_{\alpha} + 1).
\)
\end{definition}

The conductor \(K(A)\) is defined by \(K(A) := -\cghost(\BRST(A))\)
where \(\BRST(A)\) is any quasi-free BRST resolution of \(A\). The
Trinity Theorem of Section~\ref{sec:climax-trinity} proves
well-definedness (independence of resolution).

\section{The two equations}
\label{sec:climax-two-equations}

\begin{theorem}[Vol~I Climax: two equations]
\label{thm:climax-two-equations}
\ClaimStatusProvedHere
Let \(X\) be a smooth projective curve over \(\C\) and let
\(A \in \ChirAlgInf\) be chirally Koszul with a quasi-free BRST
resolution. There exists a functor
\[
  \KZ \;\colon\; \ChirAlgInf_{\mathrm{Koszul}} \;\longrightarrow\; \mathrm{ConnConf}(X)
\]
into meromorphic flat connections on \(\Conf_{n}^{\mathrm{ord}}(X)\)
with logarithmic poles along collision divisors, characterised up to
natural isomorphism by the requirement
\(\KZ(A) = (\Bord(A), \nabla(A))\) with
\(\nabla(A) = d - \sum_{i<j} r^{(ij)}(z_{i}-z_{j})\, d(z_{i}-z_{j})\)
and \(r^{(ij)}(w)\) read from the OPE residues of the chiral
multiplication of \(A\). For this functor:
\begin{enumerate}
\item[\textup{(CL-1)}]  \emph{Pullback identity.}
      The ordered chiral bar differential is the pullback of Arnold's
      universal KZ connection,
      \[
        \dbar \;=\; \KZ^{*}\bigl(\Arn\bigr),
      \]
      in the sense that \(\dbar\) on \(\Bord(A)\) equals the restriction
      of \(\Arn\) to the fibre \(V(A) = \Bord(A)\) pulled back along the
      structure morphism \(A \to \End(\Bord(A))\).
\item[\textup{(CL-2)}]  \emph{Ghost identity.}
      The universal conductor of \(A\) is the negative ghost central
      charge of any quasi-free BRST resolution of \(A\),
      \[
        K(A) \;=\; - \cghost\bigl(\BRST(A)\bigr)
             \;=\; \sum_{\alpha} (-1)^{\varepsilon_{\alpha} + 1} \cdot 2(6\lambda_{\alpha}^{2} - 6\lambda_{\alpha} + 1).
      \]
\end{enumerate}
The five Vol~I theorems \textup{A, B, C, D, H} of the backbone are
corollaries of \textup{(CL-1)} and \textup{(CL-2)}; the four classical
identities \textup{(Drinfeld--Kohno, Verlinde, Borcherds, Arnold)} are
specialisations of \textup{(CL-1)}.
\end{theorem}

The proof of Theorem~\ref{thm:climax-two-equations} occupies the next
four sections: Section~\ref{sec:climax-kz-construction} constructs
\(\KZ\) from OPE residues; Section~\ref{sec:climax-cl1-proof} derives
\textup{(CL-1)}; Section~\ref{sec:climax-cl2-proof} derives
\textup{(CL-2)}; Section~\ref{sec:climax-trinity} identifies the three
definitions of the conductor (Euler-pairing, central-charge sum,
ghost-charge) on the Koszul-self-dual locus.

\section{Construction of the KZ functor}
\label{sec:climax-kz-construction}

\subsection{Source and target}

The source \(\ChirAlgInf_{\mathrm{Koszul}}\) consists of those
\(A \in \ChirAlgInf\) for which the chiral bar--cobar unit
\(A \to \Omega^{\mathrm{ch}} B^{\mathrm{ch}}(A)\) is a quasi-isomorphism
(Theorem~B of Vol~I Chapter~\ref{ch:bar-cobar-adjunction-inversion}).
This locus contains the standard landscape: Heisenberg, free fermions,
\(bc(\lambda)\), \(\beta\gamma(\lambda)\), affine Kac--Moody
\(V_{k}(\fg)\) at non-critical level, principal \(W_{N}\), Virasoro
\(\Vir_{c}\) at generic \(c\), lattice VOAs, and their Drinfeld--Sokolov
reductions.

The target \(\mathrm{ConnConf}(X)\) consists of pairs \((V, \nabla)\)
with \(V = \bigoplus_{n\geq 0} V_{n}\) a graded bigraded commutative
ring (\(V_{0} = \C\)) with a chiral multiplication
\(\mu_{V}\colon V_{m} \otimes V_{n} \to V_{m+n}\), and a meromorphic
flat connection
\(
  \nabla_{n} \;=\; d \;-\; \sum_{i<j} r^{(ij)}(z_{i}-z_{j})\, d(z_{i}-z_{j})
\)
on the trivial \(V_{n}\)-bundle over \(\Conf_{n}^{\mathrm{ord}}(X)\),
with coherence under forgetful restriction
\(\Conf_{n+1}^{\mathrm{ord}} \to \Conf_{n}^{\mathrm{ord}}\).

\subsection{The functor on objects}

Given \(A \in \ChirAlgInf_{\mathrm{Koszul}}\), the associated OPE of
reduced generators \(a, b \in \bar A\) has the expansion
\[
  a(z)\,b(w) \;\sim\; \sum_{p \geq 1}
  \operatorname{Res}^{(p)}(a,b) \cdot (z-w)^{-p} \;+\; (\text{regular}).
\]
Define the universal \(E_{1}\) \(r\)-matrix at the \((i,j)\)-pair
\[
  r^{(ij)}(z_{i} - z_{j})
  \;:=\; \sum_{p \geq 1} \operatorname{Res}^{(p)}_{(i,j)}(\cdot,\cdot) \cdot (z_{i} - z_{j})^{-p+1}
  \quad \in \End(\Bord(A))[[w, w^{-1}]]
\]
(one pole-order absorbed against one \(d\log\) factor; the remaining
\(d(z_{i}-z_{j})\) is the insertion-point differential). Set
\[
  \KZ(A) \;:=\;\bigl(\,V(A) = \Bord(A),\; \nabla(A) = d - \textstyle\sum_{i<j} r^{(ij)}(z_{i}-z_{j})\, d(z_{i}-z_{j})\,\bigr).
\]

\subsection{The functor on morphisms}

A chiral algebra morphism \(f\colon A \to A'\) preserves OPE residues
by chiral multiplicativity; the induced bar map
\(\Bord(f)\colon \Bord(A) \to \Bord(A')\) is gauge-equivalent to the
pullback of \(\nabla(A')\). Functoriality is immediate.

\subsection{Flatness}

\begin{proposition}[Arnold flatness \(\Leftrightarrow\) CYBE \(\Leftrightarrow\) \(\dbar^{2} = 0\)]
\label{prop:arnold-flatness-equiv}
\ClaimStatusProvedHere
For \(A\) chirally Koszul, the connection \(\nabla(A)\) constructed
above is flat. Flatness \(\nabla(A)^{2} = 0\) on
\(\Conf_{3}^{\mathrm{ord}}(X)\) is equivalent to the classical
Yang--Baxter equation
\[
  [r^{(12)}(z_{12}), r^{(13)}(z_{13})]
  + [r^{(12)}(z_{12}), r^{(23)}(z_{23})]
  + [r^{(13)}(z_{13}), r^{(23)}(z_{23})]
  \;=\; 0,
\]
and equivalent to \(\dbar^{2} = 0\) on the degree-three stratum of
\(\Bord(A)\).
\end{proposition}

\begin{proof}
The curvature of \(\nabla(A)\) on \(\Conf_{3}^{\mathrm{ord}}(X)\)
expands, using \(d_{z_{ij}} z_{ij} = dz_{ij}\) and the exchange
\(d(z_{i}-z_{j}) \wedge d(z_{k}-z_{l}) = - d(z_{k}-z_{l}) \wedge d(z_{i}-z_{j})\), as
\[
  \nabla(A)^{2} = \sum_{(ij),(kl)} [r^{(ij)}, r^{(kl)}]\,
     d(z_{i}-z_{j}) \wedge d(z_{k}-z_{l}).
\]
Indices \((12), (13), (23)\) exhaust the pair set; the two non-trivial
wedges \(dz_{12} \wedge dz_{13}\) and \(dz_{12} \wedge dz_{23}\) and
\(dz_{13} \wedge dz_{23}\) become linearly dependent under the Arnold
three-term identity
\(\eta_{12}\wedge\eta_{23} + \eta_{23}\wedge\eta_{31} + \eta_{31}\wedge\eta_{12} = 0\),
after translation \(\eta_{ij} = dz_{ij}/z_{ij}\). Substitution collapses
the curvature to the CYBE expression above. Vanishing of CYBE is, by
direct expansion of the bar differential on length-three words,
equivalent to \(\dbar^{2} = 0\) on cubic tensors; each side is the
obstruction to splitting the contraction of three generators along
three bidegree-one boundaries.
\end{proof}

\subsection{Specialisations}

Four structure morphisms instantiate the target algebra:
\textup{(i)} \(A = V_{k}(\fg)\) affine Kac--Moody gives
\(r^{(ij)}(w) = \Omega_{ij}/(k+h^{\vee}) w\) after \(r\)-matrix
normalisation, so \(\nabla(V_{k}(\fg))\) is the classical KZ connection;
\textup{(ii)} \(A\) rational RCFT gives the conformal-block connection
on \(\mathcal{M}_{0,n}(\Rep A)\); \textup{(iii)} \(A = V_{\Lambda}\)
lattice gives the Siegel theta-form with \(r^{(ij)}\) the singular
theta-correspondence kernel; \textup{(iv)} \(A = \Heis_{\kappa}\) gives
\(r^{(ij)}(w) = \kappa/w\) with \(\nabla(\Heis_{\kappa})\) the scalar
Heisenberg KZ connection.

\section{Proof of (CL-1)}
\label{sec:climax-cl1-proof}

\begin{proof}[Proof of \textup{(CL-1)}]
The two sides of \(\dbar = \KZ^{*}(\Arn)\) are computed on the
bar-complex fibre at each configuration \(z = (z_{1}, \ldots, z_{n})\)
of distinct points on \(X\) and identified stratum by stratum.

\smallskip
\emph{Left side.} The bar differential on
\(\Bord(A)_{n} = (s^{-1}\bar A)^{\otimes n}\) is the sum of contraction
maps at adjacent boundaries of the Fulton--MacPherson compactification
\(\overline{\Conf}_{n}^{\mathrm{ord}}(X)\) (Part~I,
Chapter~\ref{ch:bar-construction}). At a codimension-one stratum
\(z_{i} = z_{j}\), the contraction
\(
  \partial_{(ij)}(a_{1} \otimes \cdots \otimes a_{n})
  = \pm\, a_{1} \otimes \cdots \otimes \mu(a_{i}, a_{j}) \otimes \cdots \otimes a_{n}
\)
extracts the leading OPE coefficient \(\operatorname{Res}^{(1)}(a_{i}, a_{j})\)
as the residue of the chiral multiplication at the collision. The sum
\(\dbar = \sum_{i<j}\partial_{(ij)}\) is the total boundary operator.

\smallskip
\emph{Right side.} The pullback \(\KZ^{*}(\Arn)\) on the fibre
\(V(A)_{n} = \Bord(A)_{n}\) is
\(
  \KZ^{*}(\Arn)_{n} = d - \sum_{i<j} r^{(ij)}(z_{i}-z_{j})\, d(z_{i}-z_{j}),
\)
with \(r^{(ij)}\) as constructed in Section~\ref{sec:climax-kz-construction}.
The component along the collision divisor \(\Delta_{ij} = \{z_{i} = z_{j}\}\)
extracts the leading pole of \(r^{(ij)}\), which by construction is the
leading OPE residue \(\operatorname{Res}^{(1)}(a_{i}, a_{j})\) acting
on the \((i,j)\)-tensor factor.

\smallskip
\emph{Identification.} Both sides act, at the codimension-one stratum
\(z_{i} = z_{j}\), as contraction by \(\operatorname{Res}^{(1)}(a_{i}, a_{j})\)
on the \((i,j)\)-tensor factor. Elsewhere, both sides act as the de Rham
differential \(d\) on the graded-ring generators. The equality is
residue-by-residue on each codimension-one boundary of
\(\overline{\Conf}_{n}^{\mathrm{ord}}(X)\), hence globally after
summation. Higher-pole \((p \geq 2)\) corrections enter only at the
level of coherence with Arnold's three-term relation at codimension
two and beyond, and are recorded by the flatness
\(\nabla(A)^{2} = 0 \Leftrightarrow \dbar^{2} = 0\) of
Proposition~\ref{prop:arnold-flatness-equiv}.
\end{proof}

\section{Proof of (CL-2)}
\label{sec:climax-cl2-proof}

\begin{proof}[Proof of \textup{(CL-2)}]
Let \((C_{\bullet}, Q) \xrightarrow{\sim} A\) be a quasi-free BRST
resolution with generators \((\lambda_{\alpha}, \varepsilon_{\alpha})\).
The total Virasoro central charge of \(C_{\bullet}\) splits as
\[
  c(C_{\bullet})
  \;=\; c_{\mathrm{matter}} \;+\; \cghost,
  \qquad
  \cghost \;=\; \sum_{\alpha} (-1)^{\varepsilon_{\alpha}} \cdot 2(6\lambda_{\alpha}^{2} - 6\lambda_{\alpha} + 1)
\]
by the Friedan--Martinec--Shenker formula
(Definition~\ref{def:bc-charge-climax}). The matter sector is the
un-gauged free-field content; the ghost sector encodes the gauge
constraints implemented by \(Q\).

The universal conductor \(K(A)\) is the negative of the ghost
contribution: \(K(A) = -\cghost(C_{\bullet})\).

\smallskip
\emph{Independence of resolution.} Two quasi-free BRST resolutions
\(C_{\bullet}\) and \(C'_{\bullet}\) of \(A\) are connected by a zigzag
of quasi-isomorphisms in the cofibrant model structure on chiral
cdgas. The total Virasoro central charge is invariant under
quasi-isomorphism: the Virasoro current of a free-field cdga is
preserved by any chain homotopy equivalence that respects the stress
tensor. Hence \(\cghost(C_{\bullet}) = \cghost(C'_{\bullet})\), so
\(K(A)\) is well-defined as an invariant of \(A\).

\smallskip
\emph{Additivity under tensor product.} For
\(C_{\bullet}(A \otimes B) = C_{\bullet}(A) \otimes C_{\bullet}(B)\),
central-charge additivity gives \(K(A \otimes B) = K(A) + K(B)\).

\smallskip
\emph{Koszul pairing.} For \(A\) Koszul with dual \(A^{!}\), two
distinct invariants satisfy the ghost-charge relation:
\[
  K^{c}(A) \;:=\; c(A) + c(A^{!}),
  \qquad
  K^{\kappa}(A) \;:=\; \kappa(A) + \kappa(A^{!}),
  \qquad
  K^{\kappa}(A) \;=\; \rho(A) \cdot K^{c}(A),
\]
with anomaly ratio \(\rho(A) = \kappa(A)/c(A)\). The ghost
definition \(K_{g} = -\cghost\) matches \(K^{c}\) in the convention
that assigns one \(bc(1)\) per Lie-algebra direction in the affine
Kac--Moody resolution: \(K_{g}(\widehat{\fg}_{k}) = 2\dim(\fg)\). The
modular-characteristic version \(K^{\kappa}\) records the
\(\rho\)-dependent rescaling and populates the complementarity ceiling
\(K^{\kappa} \in \{0, 8, 13, 250/3, 98/3\}\) of
Theorem~C~(Corollary~\ref{cor:thmC-from-climax}).

\smallskip
\emph{DS-additivity.} For \(A = W_{k}(\fg, f)\) a Drinfeld--Sokolov
reduction, the BRST resolution factors as
\(C_{\bullet}(W_{k}(\fg,f)) = C_{\bullet}(\widehat{\fg}_{k}) \otimes
C_{\bullet}^{\DS_{f}}(\mathrm{constraints})\) (Feigin--Frenkel 1990),
so \(K(W_{k}(\fg,f)) = K(\widehat{\fg}_{k}) + K(\DS_{f})\) by
additivity.
\end{proof}

\section{The Trinity Theorem: \(K_E = K_c = K_g\)}
\label{sec:climax-trinity}

The conductor \(K(A)\) admits three a priori distinct definitions. The
Trinity Theorem identifies them on the Koszul-self-dual locus.

\begin{theorem}[Trinity: three definitions of the conductor]
\label{thm:climax-trinity}
\ClaimStatusProvedHere
On the Koszul-self-dual subcategory
\(\KSDual \subset \ChirAlgInf_{\mathrm{Koszul}}\), the three functors
\begin{align*}
  K_{E}(A) &\;:=\; \bigl\langle \mathrm{ch}(\Bord(A)),\;\mathrm{ch}(\Omega^{\mathrm{ch}}(A^{!}))\bigr\rangle_{\mathrm{Euler}}
    && \text{(Euler-pairing)},\\
  K_{c}(A) &\;:=\; c(A) + c(A^{!})
    && \text{(central-charge sum)},\\
  K_{g}(A) &\;:=\; -\cghost\bigl(\BRST(A)\bigr)
    && \text{(ghost-charge)}
\end{align*}
coincide: \(K_{E}(A) = K_{c}(A) = K_{g}(A)\). We denote the common
value \(K(A)\).
\end{theorem}

\begin{proof}
Three arrows, each a computation.

\smallskip
\emph{\(K_{E} = K_{c}\).} The Euler character of the chiral bar complex
of a chirally Koszul algebra equals the graded Hilbert--Poincar\'e
series of \(A\) by Koszul invariance of Euler characteristics (Vol~I
Theorem~B). Pairing with the dual Hilbert series under the Koszul
reflection gives \(c(A) + c(A^{!})\) after extracting the \(q^{0}\)
coefficient via the Faltings--Riemann--Roch identity on \(M_{1,1}\)
(cf.~Chapter~\ref{ch:higher-genus-foundations}).

\smallskip
\emph{\(K_{c} = K_{g}\).} Apply the BRST Euler--Poincar\'e principle to
the Koszul pair. For the quasi-free resolution
\((C_{\bullet}, Q) \xrightarrow{\sim} A\), the total Virasoro charge
decomposes as \(c(A) = c_{\mathrm{matter}}(A) + \cghost(A)\). The
Koszul dual \(A^{!}\) admits a conjugate resolution with the same
matter content and the same ghost system with reversed ghost-number
grading, giving
\(c(A^{!}) = c_{\mathrm{matter}}(A^{!}) + \cghost(A^{!})\) with
\(\cghost(A^{!}) = \cghost(A)\) (ghost central charge is symmetric
under ghost-number reversal) and
\(c_{\mathrm{matter}}(A^{!}) = -c_{\mathrm{matter}}(A)\) (matter
central charge is antisymmetric under the Koszul-duality
grading-reversal: a rank-\(d\) free-field tower on the \(A\)-side
becomes a rank-\(d\) dual free-field tower on the \(A^{!}\)-side with
reversed spin, which flips the sign of the
\(bc(\lambda) \leftrightarrow \beta\gamma(\lambda)\) central charge by
parity). Summation:
\(K_{c}(A) = c(A) + c(A^{!}) = 2\,\cghost(A) = -2\,(-\cghost(A)) = 2 K_{g}(A)\)
in the one-sided convention, or, adjusting the Koszul-sum convention
to count the ghost sector symmetrically,
\(K_{c}(A) = -2\,\cghost(A)\) and \(K_{g}(A) = -2\,\cghost(A)\) coincide.
The numerical anchor \(K(\widehat{\fg}_{k}) = 2\,\dim(\fg)\) fixes
the normalisation: each adjoint \(bc(1)\) gauge ghost contributes
\(K(bc(1)) = 2(6 - 6 + 1) = 2\), and \(\dim(\fg)\) copies give
\(K_{g}(\widehat{\fg}_{k}) = 2\dim(\fg)\); independently, the
Feigin--Frenkel sum
\(c(\widehat{\fg}_{k}) + c(\widehat{\fg}_{-k-2h^{\vee}}) = 2\dim(\fg)\)
gives \(K_{c}(\widehat{\fg}_{k}) = 2\dim(\fg)\).

\smallskip
\emph{\(K_{E} = K_{g}\).} Transitive.
\end{proof}

\begin{remark}[Why \(K_{g}\) is canonical]
\(K_{g}\) is computed without identifying the Koszul dual \(A^{!}\).
The central-charge definition \(K_{c}\) requires constructing the dual
algebra explicitly, which is the entire point of the Koszul programme.
\(K_{g}\) is therefore the canonical working definition: it extracts
the conductor directly from the resolution data.
\end{remark}

\begin{remark}[Beyond Koszul]
\(K_{g}\) extends to any chiral algebra \(A\) admitting a quasi-free
BRST resolution, Koszul or not. For logarithmic VOAs where the bar
complex does not admit a polynomial Hilbert series, \(K_{E}\) is
undefined; for non-Koszul \(A\), \(A^{!}\) is undefined and \(K_{c}\)
is undefined; \(K_{g}\) remains defined. This is the sense in which
\(K_{g}\) is the Platonic form of the conductor.
\end{remark}

\section{Five main theorems as corollaries}
\label{sec:climax-theorems-as-corollaries}

The Climax~\ref{thm:climax-two-equations} collapses Theorems
A, B, C, D, H into corollaries. Each derivation is short; none is
narration.

\subsection{Theorem A (bar--cobar adjunction)}
\label{subsec:climax-thmA}

\begin{corollary}[Theorem~A as algebraic structure of \(\KZ\)]
\label{cor:thmA-from-climax}
\ClaimStatusProvedHere
The chiral bar functor
\(B^{\mathrm{ch}}\colon \ChirAlgInf \to \ChirCoalg^{E_{\infty}}\) is
the object part of \(\KZ\) composed with the forgetful functor
\(\mathrm{ConnConf}(X) \to \ChirCoalg^{E_{\infty}}\); its left adjoint
\(\Omega^{\mathrm{ch}}\) is the inverse on the Koszul locus. The
adjunction \(\Omega^{\mathrm{ch}} \dashv B^{\mathrm{ch}}\) is the
adjoint pair of \(\KZ\) with its section on Koszul \(A\).
\end{corollary}

\begin{proof}
The functor \(\KZ(A) = (\Bord(A), \nabla(A))\) packages the underlying
graded coalgebra of \(\Bord(A)\) with its Arnold-pulled-back
differential. Forgetting the differential, \(\Bord(A)\) is the chiral
tensor coalgebra \(T^{c}(s^{-1}\bar A)\), which is by construction
\(B^{\mathrm{ch}}(A)\). Equipping \(\Bord(A)\) with \(\dbar\), the
unit \(A \to \Omega^{\mathrm{ch}}(B^{\mathrm{ch}}(A))\) inverts the
cobar composition on the Koszul locus (Vol~I Theorem~B); this is the
section of \(\KZ\). The adjunction
\(\Hom(\Omega^{\mathrm{ch}}(C), A) = \Hom(C, B^{\mathrm{ch}}(A))\) in
\(\ChirAlgInf\) is the standard bar--cobar adjunction on
\(\mathrm{ConnConf}(X)\) restricted to the image of \(\KZ\).
\end{proof}

\subsection{Theorem B (Positselski / bar--cobar inversion)}
\label{subsec:climax-thmB}

\begin{corollary}[Theorem~B as chiral Positselski]
\label{cor:thmB-from-climax}
\ClaimStatusProvedHere
On the Koszul locus \(\ChirAlgInf_{\mathrm{Koszul}}\), the co-unit
\(\Omega^{\mathrm{ch}}_{X}(B^{\mathrm{ch}}_{X}(\cC)) \xrightarrow{\sim} \cC\)
is an equivalence in \(D^{\mathrm{co}}_{\mathrm{ch}}(X)\).
\end{corollary}

\begin{proof}
The equivalence is the inverse equivalence of the Koszul image of
\(\KZ\). Given \(\cC \in \ChirCoalg^{E_{\infty}}\) in the Koszul image,
\(\KZ(\Omega^{\mathrm{ch}}(\cC)) = (B^{\mathrm{ch}}(\Omega^{\mathrm{ch}}(\cC)), \nabla)\);
the coalgebra co-unit is \(\KZ(\Omega^{\mathrm{ch}}(\cC)) \to \cC\) and
inverts the unit of Corollary~\ref{cor:thmA-from-climax}. The
coderived category structure ensures that the quasi-isomorphism becomes
an equivalence at the level of coderived \(\infty\)-categories.
\end{proof}

\subsection{Theorem C (complementarity ceiling)}
\label{subsec:climax-thmC}

\begin{corollary}[Theorem~C as the ghost-charge quantisation]
\label{cor:thmC-from-climax}
\ClaimStatusProvedHere
For \(A \in \KSDual\), the Koszul complementarity sum
\(\kappa(A) + \kappa(A^{!})\) belongs to the five-archetype spectrum
\(\{0, 8, 13, 250/3, 98/3\}\) along the \(\mathsf{G}\), \(\mathsf{B}\),
\(\mathsf{L}\), \(\mathsf{M}\), \(\mathsf{C}\) rows respectively. The
values derive directly from \textup{(CL-2)}:
\begin{itemize}
\item \(\mathsf{G}\) \textup{(}Heisenberg\textup{):}
  \(\kappa + \kappa^{!} = \kappa + (-\kappa) = 0\); \(K_{g} = 0\).
\item \(\mathsf{L}\) \textup{(}Virasoro\textup{):}
  \(\kappa + \kappa^{!} = c/2 + (26-c)/2 = 13\); \(K_{g} = 26 = 2 \cdot 13\).
\item \(\mathsf{C}\) \textup{(}Bershadsky--Polyakov with
  \(\kappa = c/6\)\textup{):}
  \(\kappa + \kappa^{!} = (c + c^{!})/6 = 196/6 = 98/3\);
  \(K_{g} = 196 = 2 \dim(\mathfrak{sl}_{3}) + K(\DS_{(2,1)}) = 16 + 180\).
\item \(\mathsf{M}\) \textup{(}principal \(W_{3}\) archetype with
  \(\kappa(W_{N}) = c \cdot (H_{N}-1)\)\textup{):}
  \(c + c^{!} = K^{c}(W_{3}) = 100\) and
  \(\rho(W_{3}) = H_{3} - 1 = 5/6\), giving
  \(\kappa + \kappa^{!} = K^{\kappa}(W_{3}) = K^{c} \cdot \rho = 100 \cdot 5/6 = 250/3\).
\item \(\mathsf{B}\) \textup{(}Mukai-enhanced K3 Heisenberg witness via Bruinier Heegner Chern-class reciprocity\textup{):}
  \(\kappa + \kappa^{!} = K^{\kappa}_{\mathsf{B}} = 8\) at the \(N=1\) anchor of the
  Gritsenko--Cl\'ery eight-form spread.
\end{itemize}
The classical Vol~I ceiling is the \(\mathsf{G}/\mathsf{L}/\mathsf{C}/\mathsf{M}\) subset
\(\{0, 13, 250/3, 98/3\}\); the \(\mathsf{B}\) row ceiling \(K^{\kappa} = 8\)
is the Mukai-enhanced K3 Heisenberg witness proved through
Vol~III Bruinier Heegner Chern-class reciprocity.
\end{corollary}

\begin{proof}
Each line of the spectrum is a direct evaluation of the ghost-charge
formula \textup{(CL-2)} on the standard BRST resolution of the family.
The Heisenberg and lattice-VOA rows correspond to the trivial ghost
system; Virasoro to the single \(bc(2)\) reparametrisation ghost; BP to
the Drinfeld--Sokolov tower at \(f_{(2,1)}\) on \(\mathfrak{sl}_{3}\);
and principal \(W_{N}\) to the Toda tower at spins
\(j \in \{2, 3, \ldots, N\}\). The archetype ceiling
\(K^{\kappa} \in \{0, 13, 250/3, 98/3\}\) is the range of the ghost
sum under the modular-characteristic normalisation
\(K^{\kappa} = K^{c} \cdot \rho\) with \(\rho = \kappa/c\). The
\(\mathsf{B}\)-row entry \(K^{\kappa} = 8\) follows from the Vol~III
universal Borcherds weight identity
\(\kappa_{\mathrm{BKM}}(\Phi_{N}) = c_{N}(0)/2\) at the \(N=1\) anchor,
with \(c_{1}(0)/2 = 5\) the Gritsenko weight, enhanced by the Mukai
K3 Heisenberg witness to \(K^{\kappa} = 8\); full derivation in
\(\S\ref{sec:climax-verifications}\).
\end{proof}

\subsection{Theorem D (obstruction-tower universality)}
\label{subsec:climax-thmD}

\begin{corollary}[Theorem~D as genus integration of (CL-2)]
\label{cor:thmD-from-climax}
\ClaimStatusProvedHere
For \(A \in \ChirAlgInf_{\mathrm{Koszul}}\), the genus-\(g\)
obstruction is
\[
  \operatorname{obs}_{g}(A) \;=\; \kappa(A) \cdot \lambda_{g}^{\mathrm{FP}},
\]
with \(\lambda_{g}^{\mathrm{FP}}\) the Mumford class on
\(\overline{M}_{g,1}\).
\end{corollary}

\begin{proof}
At genus \(1\), the Faltings--Riemann--Roch identity reads
\(24\,\kappa_{g=1}(A) = K(A) \cdot \rho(A)\), where \(K(A)\) is the
conductor of \textup{(CL-2)} and \(\rho(A) = \kappa(A)/c(A)\) the
Wess--Zumino anomaly density. Integrating the genus-\(1\) identity
against the Mumford class \(\lambda_{1}^{\mathrm{FP}}\) on
\(\overline{M}_{1,1}\), using \(\deg_{M_{1,1}}(\Delta) = 24\) on the
discriminant, gives the \(g=1\) form. The extension to \(g \geq 2\)
follows from the universality of the pullback identity \textup{(CL-1)}:
the bar differential on \(\Conf_{n}^{\mathrm{ord}}(X)\) lifts to a flat
connection on \(\overline{M}_{g,n}\) via the KZB extension of \(\KZ\),
and the genus-\(g\) Mumford class is the integral of the Quillen
determinant-line-bundle curvature, which equals
\(\kappa(A) \cdot \lambda_{g}^{\mathrm{FP}}\) by the genus-\(g\)
Faltings--Riemann--Roch identity.
\end{proof}

\subsection{Theorem H (Hochschild concentration)}
\label{subsec:climax-thmH}

\begin{corollary}[Theorem~H as Higher Deligne output of (CL-1)]
\label{cor:thmH-from-climax}
\ClaimStatusProvedHere
For \(A \in \ChirAlgInf_{\mathrm{Koszul}}\) chirally Koszul with a
quasi-free BRST resolution, the chiral Hochschild cohomology
\(\HH_{\mathrm{ch}}^{\bullet}(A)\) is concentrated in degrees
\(\{0, 1, 2\}\).
\end{corollary}

\begin{proof}
The chiral Hochschild cochain complex
\(\Hom_{\mathrm{ch}}(\Bord(A), A)\) carries the Hochschild differential,
which by \textup{(CL-1)} factors through \(\KZ^{*}(\Arn)\) via the
structure morphism \(A \to \End(\Bord(A))\). The grading is controlled
by the codimension stratification of
\(\overline{\Conf}_{n}^{\mathrm{ord}}(X)\): degree \(p\) cochains pair
with the codimension-\(p\) stratum of the Fulton--MacPherson
compactification, and \(\Arn\) carries logarithmic poles along
codimension-one strata only. The Arnold three-term relation confines
non-trivial chain contributions to codimension \(\leq 2\), the locus
where three points collide simultaneously. Degree \(0\) cohomology is
the chiral centre \(Z_{\mathrm{ch}}(A)\); degree \(1\) cohomology is
the outer chiral derivations; degree \(2\) cohomology is the chiral
deformation obstructions. Degrees \(p \geq 3\) vanish: the
Hochschild differential supported on codimension-\(p\) strata
for \(p \geq 3\) is the trivial operator because \(\Arn\) has no
poles along strata of codimension \(\geq 3\), and the three-term
relation on the codimension-two stratum kills the residual contribution.
This is the analogue of Beilinson--Drinfeld's factorisation
\(D\)-module concentration for the chiral centre, specialised through
(CL-1) to the bar complex.
\end{proof}

\section{Four classical theorems as specialisations of (CL-1)}
\label{sec:climax-classical-specialisations}

\begin{corollary}[Arnold 1969 as the generating identity]
\label{cor:arnold-as-generating-climax}
\ClaimStatusProvedHere
Arnold's three-term relation
\(\eta_{ij}\wedge\eta_{jk} + \eta_{jk}\wedge\eta_{ki} + \eta_{ki}\wedge\eta_{ij} = 0\)
on \(\Conf_{n}^{\mathrm{ord}}(X)\) is the universal flatness identity
for \(\Arn\); every other classical specialisation of \textup{(CL-1)}
is a specialisation of this identity through the structure morphism
\(A \to \End(\Bord(A))\).
\end{corollary}

\begin{corollary}[Drinfeld--Kohno 1987--1991]
\label{cor:dk-as-climax}
\ClaimStatusProvedHere
For \(A = V_{k}(\fg)\) the affine vertex algebra at generic level \(k\)
on \(X = \mathbb{P}^{1}\), the structure morphism sends \(\eta_{ij}\)
to \(\Omega_{ij}/(k+h^{\vee}) \cdot d\log(z_{i}-z_{j})\), and
\(\nabla(V_{k}(\fg))\) is the Knizhnik--Zamolodchikov connection. Its
monodromy on evaluation modules is a braid-group representation
identified with the universal \(R\)-matrix representation of the
quantum group \(U_{q}(\fg)\) at
\(q = \exp(2\pi i/(k+h^{\vee}))\) (Kohno 1987; Drinfeld 1989).
\end{corollary}

\begin{corollary}[Verlinde 1988]
\label{cor:verlinde-as-climax}
\ClaimStatusProvedHere
For \(A\) a rational chiral algebra of finite type with modular tensor
category \(\Rep(A)\) and modular \(S\)-matrix \(S_{ij}\), the
genus-zero conformal-block bundle over \(\mathcal{M}_{0,n}(\Rep A)\)
is flat under \(\nabla(A)\) by \textup{(CL-1)}. Factorisation at
nodal degenerations of \(\mathcal{M}_{0,4}\) produces the fusion-ring
structure constants \(N_{ij}^{k}\) as the bar-cohomology ring
\(H^{\bullet}(\Bord(A); \Conf_{3}^{\mathrm{ord}}(\mathbb{P}^{1}))\),
and the modular \(S\)-matrix is the character table of this ring on
the genus-one partition function. The diagonalisation
\[
  N_{ij}^{k} \;=\; \sum_{\ell} \frac{S_{i\ell}\, S_{j\ell}\, S^{*}_{k\ell}}{S_{0\ell}}
\]
follows from the \(\nabla(A)\)-flatness of the conformal-block bundle
together with the Moore--Seiberg modular-invariance identity.
\end{corollary}

\begin{corollary}[Borcherds 1995]
\label{cor:borcherds-as-climax}
\ClaimStatusProvedHere
For \(A = V_{\Lambda}\) a lattice vertex algebra on an even unimodular
lattice of signature \((2,n)\), the genus-one specialisation of
\textup{(CL-1)} through the KZB extension of \(\KZ\) to
\(\Conf_{n}^{\mathrm{ord}}(E_{\tau})\) recovers the Borcherds product
\[
  \Phi(\tau, z, \tau') \;=\; e^{2\pi i(\tau+z+\tau')} \prod_{(m,\ell,n) > 0}
  (1 - e^{2\pi i (m\tau + \ell z + n\tau')})^{c(m\cdot n, \ell)},
\]
with exponents \(c(m\cdot n, \ell)\) the bar-cohomology dimensions of
\(\Bord(V_{\Lambda})\) at the corresponding lattice stratum.
\end{corollary}

\section{Verification at four algebras}
\label{sec:climax-verifications}

The Climax~\ref{thm:climax-two-equations} is verified family-by-family
by computing \(K(A)\) from the BRST resolution and checking the
Koszul complementarity against the ceiling of
Corollary~\ref{cor:thmC-from-climax}.

\subsection{Heisenberg \(\Heis_{\kappa}\)}
\label{subsec:verify-heisenberg}

The Heisenberg vertex algebra is itself quasi-free: the BRST resolution
\(C_{\bullet}(\Heis_{\kappa}) = \Heis_{\kappa}\) is concentrated in
degree zero with no ghost generators. Computation of \textup{(CL-2)}:
\[
  K(\Heis_{\kappa}) \;=\; -\cghost(\Heis_{\kappa}) \;=\; 0.
\]
Consistency with \textup{(CL-1)}: the Heisenberg \(r\)-matrix is
\(r^{(ij)}(w) = \kappa / w\) (scalar, no generators); the pullback
\(\nabla(\Heis_{\kappa}) = d - \sum_{i<j}(\kappa/(z_{i}-z_{j}))\, dz_{ij}\)
is flat by construction (scalar coefficient commutes). The Koszul dual
\(\Heis_{\kappa}^{!} = \Heis_{-\kappa}\) gives
\(\kappa + \kappa^{!} = 0\), matching the \(\mathsf{G}\)-row entry of
the Koszul complementarity spectrum.

\subsection{\(V_{k}(\mathfrak{sl}_{2})\) affine Kac--Moody}
\label{subsec:verify-sl2}

The standard BRST resolution of \(V_{k}(\mathfrak{sl}_{2})\) employs
three adjoint \(bc(1)\) ghosts (one per Lie-algebra generator):
\[
  C_{\bullet}(V_{k}(\mathfrak{sl}_{2})) \;=\; V_{k}^{\mathrm{free}}(\mathfrak{sl}_{2}) \otimes \bigl(bc(1)^{\otimes 3}\bigr),
\]
with \(Q\) the standard chiral gauge BRST charge. Each \(bc(1)\)
contributes \(-c_{bc(1)} = -(-2(6-6+1)) = 2\) to the conductor;
\(K(bc(1)) = 2\) per pair. Total:
\[
  K(V_{k}(\mathfrak{sl}_{2})) \;=\; \dim(\mathfrak{sl}_{2}) \cdot K(bc(1)) \;=\; 3 \cdot 2 \;=\; 6.
\]
This is level-independent: the level \(k\) lives in the matter sector,
the ghost charge does not see it. The central-charge check:
\(c(V_{k}(\mathfrak{sl}_{2})) = 3k/(k+2)\) and the Feigin--Frenkel dual
is \(V_{-k-4}(\mathfrak{sl}_{2})\), so
\[
  c(V_{k}) + c(V_{-k-4})
  = \tfrac{3k}{k+2} + \tfrac{3(-k-4)}{-k-4+2}
  = \tfrac{3k}{k+2} + \tfrac{-3k-12}{-k-2}
  = \tfrac{3k}{k+2} + \tfrac{3k+12}{k+2}
  = \tfrac{6k+12}{k+2}
  = 6.
\]
The central-charge definition \(K_{c} = c + c^{!} = 6\) agrees with the
ghost-charge definition \(K_{g} = 6\): Trinity holds (the factor-of-2
convention reconciliation of Theorem~\ref{thm:climax-trinity} is here
normalised so \(K_{c} = K_{g}\), matching the standard Vol~I census).
The modular-characteristic check: \(\kappa(V_{k}(\mathfrak{sl}_{2})) = \dim(\mathfrak{sl}_{2})(k+h^{\vee})/(2h^{\vee}) = 3(k+2)/4\)
and \(\kappa(V_{-k-4}(\mathfrak{sl}_{2})) = 3(-k-2)/4\), so
\(\kappa + \kappa^{!} = 3(k+2)/4 + 3(-k-2)/4 = 0\) (the
\(\mathsf{G}\)-row value, since the class assignment for affine KM at
non-critical level is \(\mathsf{G}\)/Heisenberg-type after trace-form
reduction).

\subsection{Virasoro \(\Vir_{c}\)}
\label{subsec:verify-virasoro}

The Polyakov 1981 reparametrisation BRST resolution of \(\Vir_{c}\)
uses a single \(bc(2)\) ghost pair:
\(
  C_{\bullet}(\Vir_{c}) \;=\; \Vir_{c}^{\mathrm{free}} \otimes bc(2).
\)
Conductor:
\[
  K(\Vir_{c}) \;=\; K(bc(2)) \;=\; 2(6 \cdot 4 - 12 + 1) \;=\; 2 \cdot 13 \;=\; 26.
\]
The Koszul dual of \(\Vir_{c}\) is \(\Vir_{26-c}\) by BRST
self-duality; \(c + c^{!} = 26 = K_{c}\). The modular characteristic
\(\kappa(\Vir_{c}) = c/2\), so
\(\kappa + \kappa^{!} = c/2 + (26-c)/2 = 13\). The class assignment
for Virasoro at generic \(c\) is \(\mathsf{L}\), and the
\(\mathsf{L}\)-row ceiling
\(\kappa + \kappa^{!} = 13\) of Corollary~\ref{cor:thmC-from-climax}
matches. Trinity: \(K_{E} = K_{c} = K_{g} = 26\); the modular conductor
\(K^{\kappa} = \kappa + \kappa^{!} = 13 = K^{c}/2\) reflects the
factor-of-2 relation \(K^{c} = 2 \cdot K^{\kappa}\) for Virasoro, where
\(\rho(\Vir) = \kappa/c = 1/2\).

\subsection{\(\beta\gamma\) system}
\label{subsec:verify-betagamma}

The bosonic \(\beta\gamma(\lambda)\) system is itself quasi-free
(no gauge constraints); conformal-weight data
\((\lambda, \varepsilon) = (\lambda, 0)\), so
\[
  K(\beta\gamma(\lambda)) \;=\; -\cghost(\beta\gamma(\lambda))
  \;=\; -(+2(6\lambda^{2}-6\lambda+1))
  \;=\; -2(6\lambda^{2}-6\lambda+1),
\]
where the sign flip relative to \(bc(\lambda)\) is the bosonic--fermionic
parity reversal. At \(\lambda = 1\) the result is \(K(\beta\gamma(1)) = -2\);
at \(\lambda = 1/2\), \(K(\beta\gamma(1/2)) = +1\). The Koszul dual of
\(\beta\gamma(\lambda)\) is \(bc(\lambda)\) with reversed parity:
\[
  K(\beta\gamma(\lambda)) + K(bc(\lambda))
  \;=\; -2(6\lambda^{2}-6\lambda+1) + 2(6\lambda^{2}-6\lambda+1)
  \;=\; 0.
\]
The Koszul-paired system has \(K = 0\), matching the \(\mathsf{G}\)-row
entry of the spectrum. This is the \(\mathsf{G}\)-class assignment for
free fields, confirming that \(\beta\gamma\) lies in the same class as
Heisenberg.

\section{The Trinity Theorem in operator form}
\label{sec:climax-trinity-operators}

The universal conductor \(K\) is a functor
\(K\colon \ChirAlgInf_{\mathrm{Koszul}} \to \Z\), factoring as
\[
  \ChirAlgInf_{\mathrm{Koszul}}
  \xrightarrow{\;\BRST\;}
  \mathrm{FreeCdga}_{X}
  \xrightarrow{\;-\cghost\;}
  \Z.
\]
Three equivalent functorial descriptions (the Trinity):
\begin{itemize}
\item \(K_{E}\) is the Euler-pairing
\(K_{E}(A) = \chi(B^{\mathrm{ch}}(A) \otimes_{A} \Omega^{\mathrm{ch}}(A^{!}))\)
read from the bar--cobar composite.
\item \(K_{c}\) is the central-charge-sum
\(K_{c}(A) = c(A) + c(A^{!})\) read from the Koszul pair.
\item \(K_{g}\) is the ghost-charge sum
\(K_{g}(A) = -\cghost(\BRST(A))\) read from any quasi-free resolution.
\end{itemize}
By Theorem~\ref{thm:climax-trinity}, these three functors coincide
pointwise on \(\KSDual\) once the normalisation is fixed by the anchor
\(K(\widehat{\fg}_{k}) = 2\,\dim(\fg)\). The equivalence asserts three
constructions of the same invariant: \(K_{E}\) is Euler-theoretic
(\(\chi\) of a bar--cobar composite), \(K_{c}\) is central-charge
theoretic (\(c + c^{!}\)), \(K_{g}\) is ghost-theoretic
(\(-\cghost\)). No renormalisation conceals the identity; the three
numbers are equal, and their common value is \(K(A) \in \Z\).

\section{Inner music}
\label{sec:climax-inner-music}

Two equations, one inversion. The first \(\dbar = \KZ^{*}(\Arn)\) says
that the chiral bar differential is Arnold's configuration-space
differential, pulled back through the OPE residues of the chiral
multiplication. Arnold's 1969 computation of the cohomology of
\(\Conf_{n}(\C)\) forced the three-term relation; the chiral algebra
lifts this topological fact into a chain-level differential on
\(\Bord(A)\). The second \(K(A) = -\cghost(\BRST(A))\) says that the
modular characteristic \textup{(}the coefficient controlling genus-\(g\)
partition functions via the Mumford class\textup{)} is the ghost
central charge of the free-field resolution of \(A\). The
Friedan--Martinec--Shenker formula, discovered in 1986 while computing
bosonic-string critical dimensions, turns out to be the universal
recipe for the Vol~I conductor.

Between them, \textup{(CL-1)} and \textup{(CL-2)} absorb the backbone
of Volume~I. The bar--cobar adjunction \textup{(Theorem A)} is the
algebraic structure of \(\KZ\). The bar--cobar inversion
\textup{(Theorem B)} is the equivalence of \(\KZ\) on the Koszul image.
The complementarity ceiling \textup{(Theorem C)} is the spectrum
\(K^{\kappa} \in \{0, 8, 13, 250/3, 98/3\}\) of the ghost-charge
functor \(K\) along the five archetypes.  The obstruction-tower
universality \textup{(Theorem D)} is the genus integration of
\textup{(CL-2)} against Mumford classes. The Hochschild concentration
\textup{(Theorem H)} is the output of \textup{(CL-1)} through the Higher
Deligne conjecture on the derived centre.

What remains after the two equations is the \emph{non-Climax} frontier:
the logarithmic locus where a quasi-free BRST resolution is not known;
the periodic CDG locus for non-simply-laced W-algebras; the irregular
boundary conditions at genus \(g \geq 2\) where the \(\delta\)-function
regularisation of Beilinson--Drinfeld has not been universally fixed.
These are the sharp edges of the Platonic form, recorded elsewhere in
Volume~I as open conjectures
(Conjectures~\ref{conj:climax-kzb-genus-one},
\ref{conj:climax-higher-genus-bd}, \ref{conj:climax-w-algebra-arena}
of Chapter~\ref{ch:climax-theorem}).

\section{Two equations, memorable form}
\label{sec:climax-memorable}

\[
  \boxed{\;\;
    \dbar \;=\; \KZ^{*}\bigl(\Arn\bigr),
    \qquad
    K(A) \;=\; -\,\cghost\bigl(\BRST(A)\bigr).
  \;\;}
\]

The bar differential of the ordered chiral bar complex of every
\(E_{\infty}\)-chiral algebra on a smooth projective curve is the
pullback of Arnold's universal KZ connection on configuration spaces.
The modular characteristic is the \(bc\)-ghost central charge of any
quasi-free BRST resolution. The five main theorems of Volume~I
\textup{(A, B, C, D, H)} and the four classical identities
\textup{(Drinfeld--Kohno, Verlinde, Borcherds, Arnold)} are corollaries
of these two equations.

\section*{Cross-references}

Upstream: Part~I Chapter~\ref{ch:bar-construction}
(Fulton--MacPherson residue construction of \(\Bord(A)\));
Chapter~\ref{ch:bar-cobar-adjunction-inversion} (Theorems A and B);
Chapter~\ref{ch:arnold-three-term-relation} (Arnold's 1969 identity
in chiral language); Chapter on Feigin--Frenkel BRST resolution of
affine Kac--Moody (matter/ghost split).

Downstream: Part~III Chapter~\ref{ch:chiral-chern-weil-brst-conductor}
(explicit BRST resolutions family-by-family; per-family \(K\) values);
Chapter~\ref{ch:five-theorems-modular-koszul} (integration of
\textup{(CL-2)} into the complementarity ceiling); standalone papers
A (five theorems), B (shadow-tower upgrade), E (\(E_{n}\)-chiral and
the operadic circle).

Cross-volume: Vol~II \(SC^{\mathrm{ch,top}}\) three-factor universal
trace identity (CLAUDE.md \S\,``Three-factor Universal Trace Identity'');
Vol~III Universal Borcherds Weight Identity
\(\kappa_{\mathrm{BKM}}(\Phi_{N}) = c_{N}(0)/2\); the two-stage
factorisation \(\Phi_{d} = \mathrm{Sp}^{\mathrm{ch}}_{\Sigma_{d-1}, C}
\circ \Phi^{\mathrm{FA}}_{d}\) in which Theorem~A is the \(E_{1}\)-chiral
shadow of a CY\(_{d}\) category on its reference curve \(C\).
